\documentclass[11pt,a4paper]{report}
\usepackage[utf8]{inputenc}
\usepackage[T1]{fontenc}
\usepackage{lmodern}
\usepackage[french]{babel}
\usepackage{geometry}
\geometry{margin=2.5cm,headheight=14pt}
\usepackage{parskip}
\usepackage{fancyhdr}
\usepackage{hyperref}
\hypersetup{hidelinks,pdftitle={Adrien Berloty (1851 - 1905)},pdfauthor={Léon Peillon}}
\providecommand{\tightlist}{\setlength{\itemsep}{0pt}\setlength{\parskip}{0pt}}
\providecommand{\releaseversion}{Édition de travail}
\setlength{\emergencystretch}{3em}
% Titres compacts pour éviter les fins de chapitre réduites à quelques lignes.
\makeatletter
\def\@makechapterhead#1{\vspace*{12pt}{\parindent\z@\raggedright
\normalfont\ifnum\c@secnumdepth>\m@ne
\Large\bfseries\@chapapp\space\thechapter\par\nobreak\vskip14pt\fi
\huge\bfseries #1\par\nobreak\vskip24pt}}
\makeatother
\pagestyle{fancy}
\fancyhf{}
\fancyhead[L]{\small Léon Peillon}
\fancyhead[R]{\small Adrien Berloty (1851 - 1905)}
\fancyfoot[C]{\thepage}
\begin{document}
\begin{titlepage}
\centering
\vspace*{1.8cm}
{\Large Léon Peillon\par}
\vspace{1.8cm}
{\Huge\bfseries Adrien Berloty\par}
\vspace{0.8cm}
{\LARGE 1851 - 1905\par}
\vfill
{\large 1er mai 1993\par}
\vspace{1cm}
{\small \releaseversion\par}
\end{titlepage}
\tableofcontents
\clearpage
% Couverture : IMGP0125.JPG ; Originaux/Couverture.jpg ; Adrien Berloty 1851 - 1905 ; Léon Peillon 1993

% Page source 01 : IMGP0126.JPG ; Originaux/Page 01.jpg

\chapter{Avant-propos et origines familiales}\label{avant-propos-et-origines-familiales}

L'ouvrage « Adrien Berloty 1851 - 1905 » publié en 1910 sous la signature du père Rosette S.J. souhaitait essentiellement, d'après l'auteur lui-même, faire revivre et fixer la physionomie du grand chrétien et de l'homme de bien que fut Adrien Berloty.

Attachant à cet effet plus d'importance à ce que disent les lettres écrites ou reçues par A.B., le père Rosette a laissé quelque peu dans l'ombre l'enchaînement quotidien des faits qui ont constitué la trame de cette vie somme toute assez courte ; un certain désordre chronologique s'ajoutait à une présentation souvent très succincte laissant le lecteur insatisfait.

L'étude qui suit reprend la description de vie d'A.B. sous une forme moins enluminée et plus concrète en ajoutant aux renseignements tirés du père Rosette ceux qu'apportent plusieurs papiers de famille tels que les « Souvenirs de jeunesse de Félix Berloty », et divers testaments et actes de partage communiqués par Michelle Gindre.

\begin{center}\rule{0.5\linewidth}{0.5pt}\end{center}

Adrien Berloty (A.B.), né à Lyon le 3 Février 1851, était le quatrième enfant de François Félix Berloty (F.F.B.) notaire à Lyon qui avait épousé en Août 1839 Jeanne Françoise Marie Coste. F.F.B. était lui-même fils de Claude Denis Berloty et de Scholastique Bonaventure Poulat ; tous deux originaires des environs de Bourg-en-Bresse ; si A.B. ne connut pas son grand-père Claude Denis décédé en 1852 à Gros Bois, il ne cessa d'admirer pendant trente ans sa grand-mère Scholastique qui, jusqu'à son décès en 1883 à Gros Bois, resta une des personnalités de la famille.

La mère d'Adrien, Jeanne Françoise Marie Coste, était la fille de Benoît Coste et de Joséphine Colomb de Gast. De ses grands parents maternels A.B. ne put connaître que Joséphine sa grand-mère mais, sans qu'il en garde un souvenir précis et vivant, celle-ci étant décédée en 1855 alors qu'il n'avait que quatre ans. Son grand père Benoît Coste avait quitté ce monde bien avant sa naissance, étant décédé en 1845 à Montréal au Canada où, pour échapper à des créanciers tenaces, il s'était réfugié à la suite de mauvaises affaires. Comme, ce faisant, il avait littéralement abandonné sa famille, celle-ci lui en gardait quelque rancune et A.B. ne dut en entendre parler qu'à mots couverts ; le père Rosette d'ailleurs ne cite pas une seule fois le nom de Benoît alors qu'Adrien, comme cela se verra plus loin, a eu en fait très souvent, sur le plan religieux comme sur le plan politique, un comportement hérité tout-droit de son grand père Coste.

% Page source 02 : IMGP0127.JPG ; Originaux/Page 02.jpg

À sa naissance, en Février 1851, un frère né avant lui étant décédé en bas âge, A.B. se trouva le fils aîné tout en ayant deux sœurs plus âgées que lui, Joséphine née 10 ans avant lui et Blandine qui décéda de maladie en 1860. Denise née en 1852 et Bonaventure né en 1854 vinrent compléter par la suite la famille qui en définitive se stabilisa à quatre enfants.

A.B. naquit dans l'immeuble : 10, Place des Terreaux où F.F.B. avait à cette époque son appartement et son Étude. C'est dans cet immeuble ainsi qu'à Gros Bois, dont la construction était en voie d'achèvement qu'il vécut ses premières années. Il n'est pas impossible que F.F.B. ait aussi disposé à cette époque d'une location avec jardin située Montée des Carmes Déchaussés au-dessus de Saint-Paul, car c'est Montée des Carmes Déchaussés que Joséphine Coste est décédée en 1855.

\begin{center}\rule{0.5\linewidth}{0.5pt}\end{center}

Les premières années d'A.B. n'allèrent pas sans soucis pour ses parents. Dès 1854, Adrien, étant à Vichy avec ceux-ci qui y faisaient une cure, fut atteint d'une bronchite exigeant un traitement très énergique qu'il supporta mal, on le ramena à Gros Bois littéralement entre la vie et la mort. Il s'en tira finalement mais sa santé en resta délicate pendant toute son enfance. Ses premières années de classe dans différentes écoles élémentaires du quartier des Terreaux furent souvent interrompues par des indispositions plus ou moins graves au point que, de 1859 à 1862, ses parents jugèrent préférable de le placer dans un internat situé à la Mulatière et spécialisé dans l'enseignement des enfants délicats.

F.F.B. estimant alors qu'il fallait à son fils une meilleure ambiance de travail n'hésita pas sur les moyens et le fit entrer en Octobre 1862 pensionnaire au Collège de Mongré, à Villefranche-sur-Saône, Établissement réputé, fondé depuis peu par les Jésuites, et doté de cours de récréation spacieuses et ombragées dans un parc de 20 hectares.

Adrien n'était pas entré depuis deux mois à Mongré que sa mère Marie Coste décéda le 1er Décembre 1862 à la suite d'une courte maladie d'origine cardiaque. Il avait à peine 12 ans ; si sa sœur Joséphine, qui avait épousé en 1861 J.P. Boffard, menait déjà sa vie à part, Denise n'avait que 10 ans et Bonaventure 8 ans. Les trois derniers enfants de F.F.B. n'étaient que de jeunes enfants ayant encore besoin d'une présence maternelle. Heureusement leur grand-mère Scholastique Bonaventure était là avec tous ses moyens. Elle prit résolument ses petits enfants en charge, et, à cet effet, s'installa avec eux dans l'immeuble de la place de la Bourse dont la construction venait de s'achever et où F.F.B. avait transféré son appartement et son Étude.

% Page source 03 : IMGP0128.JPG ; Originaux/Page 03.jpg

D'après le père Rosette, Adrien fut très sensible à la disparition de sa mère mais on peut se demander si, simultanément, cela ne lui facilita pas son adaptation à la vie du collège. Il y trouva ce qu'il ne pouvait avoir place de la Bourse, des camarades de jeux dont il se fit de véritables amis, avec lesquels les récréations étaient certainement animées et avec lesquels par exemple il jouait régulièrement la comédie lors des fêtes du collège, y tenant même des rôles importants. Ces amis, il les invitait place de la Bourse pendant les vacances en recommandant bien à son père d'éloigner Denise ce jour-là pour ne pas les gêner ; la vie quotidienne méritait sans doute à ses yeux d'être organisée à sa manière.

Devenu veuf F.F.B. n'abandonna à personne d'autre l'éducation et la formation de son fils. Se tenant très au courant de la façon dont Adrien se comportait à Mongré, il lui écrivait souvent et, quand il le fallait, des lettres où l'affection savait s'allier à la fermeté paternelle indispensable ; les quelques lettres citées par le père Rosette le montrent bien ; F.F.B. ne cessa pas non plus d'organiser pour Adrien pendant les vacances d'été des voyages susceptibles de concourir à sa formation.

Sur le plan des études, les sept années passées à Mongré et dont le père Rosette parle assez longuement, ne se passèrent pas sans problèmes. Le père Rosette insiste sur le fait que, dès le début, A.B. se montra convaincu, d'une façon assez exceptionnelle pour un garçon de cet âge, des lourdes responsabilités qu'il aurait à assumer à la suite de son père et, par là même, incité à travailler et à obtenir un bon classement. Il semble qu'il ait été très vite estimé des Jésuites dont l'un d'eux, le père Richard, appréciant son caractère, s'attacha réellement à lui. Malheureusement les résultats ne furent pas toujours à la hauteur des espérances et Adrien en souffrit certainement.

Pendant les premières années à Mongré, en effet il s'était trouvé habituellement dans la bonne moyenne de la classe et il lui arriva même d'être premier en diligence, mais à partir de la classe de troisième les résultats se détériorèrent et il lui fallut annoncer à son père de très mauvaises places conséquences de crises d'indolence, voire de paresse ; ses professeurs qui s'intéressaient à lui, l'aidaient à les surmonter temporairement. On y voit apparaître un aspect du caractère d'Adrien qui reparut par la suite plusieurs fois au cours de sa vie ; devant des échecs répétés, il se laissait aller à des crises de découragement, doutant de sa capacité à devenir ce qu'il souhaitait et encore plus ce que souhaitait son père ; l'avenir lui paraissait en noir, il perdait goût au travail et était prêt à tout abandonner.

% Page source 04 : IMGP0129.JPG ; Originaux/Page 04.jpg

En dernière année avec la perspective du baccalauréat, il se trouva un jour, en Janvier, 25ème en version latine et du coup il écrivit à F.F.B. pour que celui-ci obtienne de la direction du collège de le dispenser de « colles » qui l'exaspéraient littéralement et qu'il appréhendait ; allant même plus loin il demanda à son père de le sortir de Mongré, alors que le recteur du collège en était venu à se demander de son côté s'il convenait bien de le garder. Le père Richard obtint heureusement un sursis et F.F.B. sut trouver les mots pour calmer son fils qui se remit au travail jusqu'à l'épreuve du baccalauréat abordée avec appréhension. En Juillet ce fut l'échec, durement ressenti, mais Adrien ne perdit pas courage et finalement se rattrapa à l'épreuve de repêchage en Novembre.

Pour détendre son fils pendant les vacances et lui changer les idées, F.F.B. l'engagea à participer au cours de l'été 1866 à un tour d'Espagne à pied, sac au dos, organisé par un cousin de ses parents, l'abbé Neyrat. Celui-ci, dont la mère était née Colomb de Gast, était un marcheur infatigable ; il ne craignait pas d'aller de Montbrison à Lyon à pied dans la journée et savait ce que c'est que d'apprendre à de jeunes garçons à dominer leur fatigue. Les premières journées de ce périple espagnol furent pour Adrien un véritable supplice mais, petit à petit, il s'y mit à tel point que, l'Abbé Neyrat étant tombé malade à la fin du voyage, ce fut lui qui le soigna et il s'en fit un véritable ami. L'année suivante, dans l'été 1867, l'abbé Neyrat ayant mis sur pied un voyage en Écosse, Adrien partit à nouveau avec lui et trouva probablement goût à visiter ce pays puisqu'il y retourna beaucoup plus tard. Il avait découvert quel agrément et quel intérêt représentaient les voyages à l'étranger, même s'ils sont fatigants.

En Novembre 1869, Adrien était donc bachelier. Depuis des années il savait quelle carrière l'attendait ; il n'y avait pas à hésiter et deux jours après son examen il s'inscrivit à la Faculté de Droit de Lyon pour préparer sa licence, étape indispensable pour pouvoir par la suite remplacer son père dans l'Étude notariale. En même temps afin de commencer à apprendre le métier, il entrait comme clerc à temps partiel à l'Étude. Automne 1869, hiver et printemps 1870 furent pour lui des mois où il oublia complètement les difficultés passées, goûtant les facilités de la vie d'étudiant avec ses libertés, ses distractions et la camaraderie ; il eut la possibilité de faire de la musique et notamment de pratiquer le chant, ce qu'il appréciait tout particulièrement.

% Page source 05 : IMGP0130.JPG ; Originaux/Page 05.jpg

Au printemps 1870, la guerre avec l'Allemagne changea brusquement l'atmosphère, surtout avec les désastres militaires de l'été. À la fin Octobre, malgré ces désastres (ou peut-être à cause d'eux) F.F.B. décida de se rendre à Cadix pour apprécier sur place la situation des affaires qu'il y avait montées quelques années auparavant ; pour ce voyage, qui impliquait de traverser l'Espagne du Nord au Sud, F.F.B. jugea opportun d'emmener avec lui Adrien afin que celui-ci commençât à voir de quoi il s'agissait. Ce fut le premier voyage d'affaires d'Adrien ; il n'avait pas 20 ans.

Quand ils rentrèrent à Lyon en Décembre 1870, la guerre continuait et l'appel sous les drapeaux de la classe 91 était envisagé. Attiré par un de ses amis engagé dans le régiment des zouaves de Charette à Rennes, Adrien songea à l'y rejoindre. C'était un régiment d'élite, exigeant donc une excellente santé, et à haut risque. F.F.B. calma sans doute rapidement l'ardeur de son fils, qui finalement demanda à s'engager dans le Génie. Avant de souscrire son engagement, il fit un voyage éclair en Suisse pour y revoir la grand-mère Scholastique qui, comme beaucoup de Lyonnais, s'était réfugiée dans ce pays afin d'y être à l'abri en cas d'invasion. À son retour à Lyon, l'armistice de Janvier 1871 venait d'être signé, l'engagement tomba de lui-même ; Adrien reprit avec plaisir la préparation de la licence, son travail à l'Étude et ses réunions avec ses amis, pensant avoir devant lui quelques années tranquilles d'une adolescence agréable et sans gros soucis et l'accès assuré à un notariat très classique. Mais brusquement, trois mois après, toutes ces perspectives furent bouleversées.

\chapter{La succession de François Félix Berloty}\label{la-succession-de-franuxe7ois-fuxe9lix-berloty}

Au début de Mai 1871, François Félix Berloty dut subir une petite intervention chirurgicale assez banale liée à un problème d'anthrax ; il souffrait malheureusement de diabète depuis plusieurs années et l'opération fut rapidement suivie de complications graves qui aboutirent à une issue fatale. F.F.B. décédait le 26 Mai 1871 dans sa 63ème année.

Adrien adorait son père pour lequel il avait beaucoup d'estime. Il fut sentimentalement très affecté par cette disparition ; le père Rosette le souligne en citant plusieurs lettres écrites ou reçues par A.B. à cette occasion, lettres d'un contexte assez conventionnel, comme il se doit en pareille circonstance, et qui n'évoquent pratiquement pas la situation toute nouvelle dans laquelle A.B. se trouvait brusquement plongé sans y avoir été préparé autrement que par la solide formation qu'il avait reçue.

% Page source 06 : IMGP0131.JPG ; Originaux/Page 06.jpg

Tout juste âgé de 20 ans et n'ayant pas encore eu le temps d'acquérir une grande expérience personnelle, il se trouvait, en tant que fils aîné, orphelin de père et de mère, promu au rang de chef de famille avec, plus âgée que lui sa sœur Joséphine mariée depuis 10 ans à J.P. Boffard mais aussi avec une sœur et un frère encore mineurs ; Denise 19 ans, en âge d'être mariée et Bonaventure 17 ans appelé à choisir rapidement une orientation. De ses ascendants susceptibles de lui apporter une aide et des conseils, il ne restait que la grand-mère Scholastique, toujours en très bonne santé mais quand même âgée de 85 ans. Seul, son beau frère Joanny Boffard, qui approchait de la quarantaine, avait de l'expérience ; il joua d'ailleurs pleinement la carte familiale en y occupant avec Joséphine une place très importante.

La disparition de F.F.B. n'apporta sans doute pas beaucoup de changements dans les soucis de la vie quotidienne, place de la Bourse, où Scholastique tenait, depuis la mort de sa belle-fille, le rôle de maîtresse de maison. L'ouverture de la succession de F.F.B. en souleva par contre beaucoup d'autres et de plus importants.

Il y avait d'abord l'avenir de l'Étude. N'ayant pas encore sa licence de droit, Adrien ne pouvait espérer être agréé comme notaire avant plusieurs années, il fallait donc céder l'Étude et cela n'allait pas sans risquer de la voir échapper définitivement. Ce problème dut tracasser immédiatement Adrien jusqu'au plus profond de lui-même mais aucune trace n'a été gardée de son inquiétude à ce sujet. J.P. Boffard apporta la solution en décidant de reprendre l'Étude à son compte jusqu'à sa retraite, avec l'engagement de la céder alors à son beau-frère. Il fut rapidement agréé comme notaire et l'Étude resta dans la famille. Adrien y conservait son activité de clerc et la possibilité d'y apprendre le métier de notaire dans les meilleures conditions ; accessoirement, il conservait aussi un lieu de travail immédiatement en-dessous de l'appartement que la famille continuait d'habiter. Le problème fut ainsi parfaitement et vite résolu et l'on n'entendit plus parler.

La succession proprement dite posa ensuite des problèmes d'autant plus complexes que la minorité de deux des héritiers ne simplifiait pas les choses et il y eut là de quoi occuper les pensées et les réflexions d'Adrien dès le mois de Juin 1871.

% Page source 07 : IMGP0132.JPG ; Originaux/Page 07.jpg

F.F.B. laissait une fortune immobilière très importante : deux immeubles à Lyon, le château de Gros Bois à Ouroux avec 600 Hectares, des vignes sur 25 Ha à Lantignié, 280 Ha à Vougy à côté de Roanne, une grande propriété à Montcarra près de Bourgoin, une autre moins importante à Saint-André-le Désert entre Cluny et Charolles, une propriété à Saint-Régis-du-Coin en Haute Loire provenant de la succession Coste en même temps que divers petits immeubles sis à Saint-Étienne en indivision avec d'autres héritiers Coste. Tout cela dans un rayon n'excédant pas 100 Km autour de Lyon mais, à cette époque où l'automobile n'existait pas encore, un déplacement de 100 Km aller et retour ne pouvait pas facilement se faire dans la journée. Le seul souci de la bonne gestion de tous ces biens immobiliers avait de quoi préoccuper sérieusement Adrien, même aidé de J.P. Boffard.

À cela vint s'ajouter que la succession apparut rapidement présenter, à côté d'un actif très important, un passif non négligeable, les propriétés acquises les plus récemment n'étant pas encore entièrement payées. Il allait donc falloir revendre une partie de ces propriétés. Lesquelles et dans quelles conditions ?

Il y avait aussi le testament que F.F.B. avait rédigé en 1867 et corrigé en Août 1870. Ce testament, exagérément précis dans ses détails notamment pour le partage des biens immobiliers, soulevait de nombreuses difficultés d'application accrues d'ailleurs par la présence d'enfants mineurs, et les quatre héritiers décidèrent finalement, avec la bénédiction de Scholastique, de refuser ce testament et de prendre des dispositions d'attente.

Il fut convenu de ne pas procéder de suite à un partage des biens immobiliers de la succession et de laisser dans une indivision à parts égales tous ceux que l'on pourrait garder après avoir vendu ce qui serait nécessaire pour apurer le passif, tout ceci sans toucher à ce dont Scholastique bénéficiait déjà c'est-à-dire la pleine jouissance jusqu'à son décès du château de Gros Bois et de son parc. C'était une solution sage mais Adrien allait se trouver pendant plusieurs années plongé dans la gestion d'une indivision immobilière composée de biens ruraux très dispersés tandis que, à Gros Bois même, il ne serait pas à proprement parler, seul propriétaire des lieux.

Tout cela dut beaucoup le préoccuper pendant tout l'été 1871. On connaît très peu de choses de ses réactions pendant toute cette période mais, sans doute sous le poids de la fatigue, il fut pris au début de l'automne d'une profonde crise de découragement analogue à celles qui l'avaient saisi précédemment à Mongré.

% Page source 08 : IMGP0133.JPG ; Originaux/Page 08.jpg

À cela vint s'ajouter une complication du type bronchite suffisamment tenace et sérieuse pour qu'un moment l'on craignît le pire, le père Rosette indiquant même qu'A.B. se préparait à la mort.

Ses proches le convainquirent de s'arrêter de travailler et de partir se soigner dans le midi. Il s'installa donc en fin 1871 à Nice dans la pension Saint-Étienne où il trouva l'accueil adéquat, plusieurs de ses amis se trouvant d'ailleurs à Nice à cette époque. Ce ne fut pas sans peine qu'il se résigna à abandonner son travail et ses responsabilités pour y passer deux mois. La bronchite n'était heureusement qu'une bronchite et la constitution au total assez solide d'Adrien en eut finalement raison.

En Mars 1872, il regagna Lyon et put y reprendre sa vie normale. Sa santé était bien rétablie ; il semble qu'il s'en inquiétait néanmoins de temps en temps car dans l'été 1873, il fit une saison à Cauterets en compagnie d'ailleurs du père Richard.

En Septembre 1873, Adrien obtint sans peine sa licence de droit, en présentant une thèse, dont le sujet ne nous est pas connu, mais qui, relue par lui quelques années après, lui parut, dit-il, avec le recul, parfaitement valable.

\chapter{Mariage, famille et premières acquisitions}\label{mariage-famille-et-premiuxe8res-acquisitions}

Sa sœur Denise, ayant épousé en Juin 1872 Joseph Rambaud, il ne restait plus au domicile familial, place de la Bourse, que Bonaventure qui terminait ses études et décida en Septembre 1874 d'entrer chez les Jésuites dans leur noviciat de Lons-le-Saunier. Les deux mineurs étant ainsi casés, Adrien se trouvait tout seul place de la Bourse avec sa grand-mère et ses 88 ans. Il n'avait plus qu'à se marier.

Il aborda la chose avec le sérieux qui le caractérisait, craignant toujours de ne pas faire assez bien. Le père Richard ami fidèle, l'encouragea à fonder une famille. Son choix se porta sur Anne Marie Demoustier qu'il avait eu plusieurs fois l'occasion de rencontrer.

C'était la fille d'Armand Romain Demoustier et d'Anne Marie Phelip. Elle avait un frère Elysée célibataire et trois sœurs : Hélène qui était religieuse au Cénacle, Marie Christine qui avait épousé en 1872 René Brunet-Leconte et devait décéder d'ailleurs quelques semaines après le mariage d'Adrien et d'Anne-Marie en laissant deux orphelins, et la troisième Elizabeth encore une jeune fille.

% Page source 09 : IMGP0134.JPG ; Originaux/Page 09.jpg

Romain Demoustier habitait Chemin de la Favorite à Lyon, une propriété de 3 Ha avec Parc et jardin qu'il avait acquise en 1855 et complétée en 1864. Intéressé par tout ce qui concernait l'horticulture, il en avait fait une propriété très agréable avec des fleurs superbes, des fruits excellents, et même quelques plantes exotiques. Il possédait et gérait personnellement une bonne charge d'Agent de Change sur la place de Lyon. Il était également propriétaire de plusieurs immeubles situés à l'angle de la place Carnot et de la rue de Condé, achetés en 1869, et d'un immeuble situé : 34, quai de la Charité, qu'il avait fait construire sur un terrain de 780 m2 propriété de son père.

Le mariage eut lieu le 6 Avril 1875 ; le contrat de mariage avait été signé quelques jours auparavant sous le régime de la communauté de biens réduite aux acquêts, Romain Demoustier faisant donation à sa fille, en avance d'hoirie, d'une somme de 120.000 F. assez importante pour l'époque. Le ménage n'allait pas manquer de ressources.

Romain Demoustier faisait de fréquents séjours en Suisse où il avait d'excellents amis. À l'automne 1870, craignant l'invasion allemande, il avait installé sa femme et ses enfants à Fribourg et Anne Marie garda bon souvenir de ce séjour. Cette attirance de la Suisse rejoignit curieusement celle que les Coste avaient pour le même pays depuis qu'en 1793 Isaac Coste\footnote{Le père de Benoît et donc l'arrière-grand-père d'Adrien.} avait été obligé de s'y réfugier ; l'attrait de la Suisse pour Adrien s'en trouva très vraisemblablement renforcé.

Le jeune ménage s'installa place de la Bourse au 2ème étage, au-dessus de l'Étude, dans l'appartement occupé autrefois par F.F.B. et dans lequel Scholastique ne garda pour elle que quelques pièces qui lui assuraient d'ailleurs un bon logement.

Quoiqu'on ait peu de renseignements sur la vie du jeune ménage pendant les premières années qui suivirent le mariage, tout permet de penser qu'il y avait en hiver les dîners de famille chez les uns et les autres, de fréquentes visites à la Favorite, surtout au printemps, et en été des séjours à Gros Bois. Existence peut-être un peu monotone, agrémentée de quelques voyages, comme le dit le père Rosette qui ne donne aucune indication sur leur localisation.

Depuis que l'Étude était entre les mains de son beau-frère, rien n'avait été changé aux occupations qu'Adrien y trouvait. Une fois la licence acquise, il y consacra tout son temps d'abord dans une position subalterne mais ensuite vraisemblablement comme clerc principal.

% Page source 10 : IMGP0135.JPG ; Originaux/Page 10.jpg

Il semble que très vite Adrien fit connaissance avec les communautés religieuses qui allaient tant l'occuper par la suite. C'est en effet à ce moment qu'il fut choisi comme légataire universel par Melle Marie Victor Cézard de Compigny cliente de l'Étude qui vivait dans une maison religieuse de Fourvière ; elle y décéda le 3 Décembre 1876, laissant un testament qui prévoyait toute une série de donations à faire après son décès. Adrien était donc déjà connu et estimé sur la colline ; cette charge de légataire universel demanda des interventions de sa part pendant de très nombreuses années.

Une lettre de 1877, que cite le père Rosette et qui parle des élections législatives au cours de cette même année, montre également que dès cette époque Adrien s'intéressait à tout ce que l'on pouvait craindre au sujet de la politique relative à la liberté d'action des institutions religieuses, politique qui n'en était qu'à ses débuts.

En 1877, deux ans seulement après son mariage et alors que rien n'était encore décidé pour la dévolution définitive des biens issus de la succession de F.F.B., Adrien procéda à sa première acquisition immobilière, acquisition, qui comme on le verra, fut suivie de beaucoup d'autres. En Juin 1877 en effet il acheta une propriété rurale de 54 Ha dite du Siber et située à proximité immédiate de Thizy sur les communes du Marnand et à Saint-Jean-la-Bussière, des prés, des terres, des vassibles et des bois avec logement pour le fermier et bâtiments d'exploitation mais pas de maison de maître, le tout dans un secteur montagneux et de faible rendement agricole sur une grande pente orientée à l'Ouest avec la maison de ferme à 600 m d'altitude. Avec le recul du temps, on ne voit pas bien quel pouvait être l'intérêt de cet achat. Peut-être Adrien voulut-il simplement suivre l'exemple de son père qui, jusqu'à la veille de sa mort, ne cessa d'acheter des terres agricoles. D'ailleurs, une fois entré en possession de Gros Bois, Adrien ne semble pas s'être beaucoup intéressé à Siber ; il n'y a pas trace qu'il y ait fait des investissements d'amélioration.

Le premier enfant se fit un peu attendre, aussi quand Marie vint au monde au début de 1878 fut-elle particulièrement bien accueillie. Adrien l'abandonna toutefois aussitôt quelques semaines, ayant été convoqué au mois d'Août pour une période de réserviste militaire de 30 jours qu'il fit dans la région de Digne. Le père Rosette ne dit pas dans quelle arme se fit cette période, mais les lettres qu'il cite montrent qu'Adrien en garda un souvenir mitigé avec des réactions mi-désabusées, mi-révoltées au fur et à mesure qu'il découvrait, au sens propre du terme, le comportement du troupier banal si différent du sien sur le plan religieux et sur le plan moral. Cette période de réserviste fut suivie d'une deuxième de même durée en 1880.

% Page source 11 : IMGP0136.JPG ; Originaux/Page 11.jpg

La fin de l'année 1878 vit intervenir un événement qu'Adrien attendait certainement avec impatience. Sept années s'étaient écoulées depuis le décès de F.F.B. Tandis que les familles de ses deux sœurs s'étaient accrues de nombreux enfants, lui-même était maintenant marié depuis trois ans et père de famille, son frère Bonaventure semblait bien ancré chez les Jésuites. Chacun aspirait à être chez soi et à sortir de l'indivision. Les quatre héritiers décidèrent de faire évoluer la situation en convenant de procéder à un « partage en nature » de tous les biens immobiliers qui prendrait effet le 1er Janvier 1879 mais sans que les valeurs individuelles des biens partagés soient fixées, cette dernière formalité étant remise à une date ultérieure ; celle-ci ne devait en fait intervenir qu'en 1887 plusieurs années après le décès de Scholastique.

Cette sortie d'indivision ne fut toutefois que partielle pour Adrien alors qu'elle devint totale pour Joséphine et Denise. Bonaventure ayant 23 ans était bien majeur, mais il n'avait pas encore terminé son noviciat chez les Jésuites en Angleterre. Adrien jouant pleinement son rôle de chef de famille accepta de rester en indivision avec lui encore quelques années, ce qui permettait d'éviter jusqu'à nouvel ordre toute dispersion du patrimoine immobilier de la famille. Le partage en nature décida donc qu'Adrien et Bonaventure deviendraient propriétaires en indivision du château de Gros Bois, de ses dépendances et des 130 Ha de bois attenant. Joséphine Boffard et Denise Rambaud se partageaient les quelques 800 Ha restant, situés à Ouroux, Vauxrenard, Lantignié et Vougy ; Joséphine gardait l'immeuble de la place de la Bourse ; le solde, non encore vendu des biens de la succession F.F.B. était à mettre en vente. C'était une solution de sagesse.

Ainsi au 1er Janvier 1879, alors qu'il atteignait 28 ans, Adrien se trouvait enfin à peu près chez lui à Gros Bois où Scholastique conservait néanmoins, conformément aux dernières volontés de F.F.B. la pleine jouissance du château et de ses dépendances, tandis que Bonaventure n'y venait plus. Cette consolidation de sa situation personnelle aurait pu inciter Adrien à mettre un terme, au moins provisoire, à de nouvelles acquisitions de biens immobiliers ruraux ; il n'en fut rien.

Début 1879 en effet, Adrien eut connaissance de la mise en vente d'une propriété importante qui le tenta, bien que située dans le Lot et donc à beaucoup plus de 100 Km de Lyon. La terre des Escouannes, 91 Ha, n'était pas inconnue de l'Étude Berloty. Elle avait appartenu à un client de l'Étude, Mr.~Victor Deville de Fontaines-sur-Saône qui l'avait vendue à un entrepreneur parisien Mr.~Dephieux, le 31.12.1867 devant François F. Berloty.

% Page source 12 : IMGP0137.JPG ; Originaux/Page 12.jpg

Mr.~Dephieux ayant fait de mauvaises affaires, la propriété venait d'être l'objet d'une saisie avec mise en vente devant le tribunal de Figeac. Celui-ci, en recherchant les origines de propriété, s'adressa évidemment à l'Étude et c'est probablement ainsi qu'Adrien eut connaissance de cette vente. N'étant pas encore titulaire de l'Étude, il avait la possibilité de se porter acquéreur et en Avril 1879 il alla passer plusieurs jours à Figeac pour réaliser l'opération, le tribunal lui adjugeant finalement la propriété sur une offre de 150.000 F.

La terre des Escouannes, située à cheval sur les communes de Prudhomat et de Loubressac, à côté de Saint-Céré et donc à la frontière entre les départements du Lot et de la Corrèze, était implantée pour l'essentiel sur une île située sur la rive droite du lit de la Dordogne, immédiatement en aval de l'endroit où la Cère se jette dans cette rivière. L'île est longue de 3 kilomètres avec une largeur atteignant jusqu'à 800 mètres, inondable en cas de crue, ce qui est presque toujours favorable au point de vue agricole en raison des apports en limon, située à basse altitude (125 mètres environ) et facile à cultiver parce que principalement en terrain plat ; cette propriété ne devait pas être sans intérêt. Nous ne savons pas de quoi elle était constituée au moment de l'achat mais, d'après l'acte de partage de la succession d'Adrien établi en 1907, alors qu'elle se trouvait réduite à 49 Hectares par suite de ventes partielles, elle comportait encore 18 Ha de terres, 9 Ha de prés, 10 Ha en pâtures et un cheptel de 18 bovins et 80 brebis.

Adrien vit là sans doute un placement intéressant, très différent par sa consistance de ce qu'était Gros Bois, et il ne recula pas devant les contraintes que la gestion de cette terre allait entraîner. Elle était exploitée en métayage et non en fermage ; il faudrait donc aller au moins une fois par an voir sur place ce qui s'y passait. À cette époque le voyage ne pouvait se faire qu'en train en passant par Clermont, Ussel, Tulle et Brive et avec plusieurs changements, il fallait plus de 24 heures pour faire un trajet qui à vol d'oiseau n'excédait pas 250 Km. Adrien, qui fut toute sa vie un voyageur infatigable, ne semble pas avoir souffert de la chose mais son fils Félix dans ses « Souvenirs de jeunesse » ne cache pas qu'obligé lui aussi à aller par la suite une fois par an aux Escouannes, il trouva plus agréable de faire le trajet en voiture automobile dès que cela fut possible.

L'année 1879, décidément bien remplie pour Adrien, vit ensuite la décision prise en Octobre par J.P. Boffard, dont la santé était ébranlée, de prendre sa retraite en cédant l'Étude à son beau-frère comme cela avait été convenu neuf ans plus tôt.

% Page source 13 : IMGP0138.JPG ; Originaux/Page 13.jpg

Les démarches à faire par Adrien en vue d'obtenir son agrément comme notaire seraient à faire au printemps 1880 de façon que la désignation officielle intervienne dans l'été, et sans attendre celle-ci il racheta l'Étude pour 300.000 F (dont 32.700 F pour les meubles) le 31 Mai 1880.

\chapter{Notaire, propriétaire et homme public}\label{notaire-propriuxe9taire-et-homme-public}

Adrien prêta serment devant le Procureur de la République le 28 Août 1880 et le 30 Août il prit officiellement ses fonctions de notaire.

Le cadre de vie familiale et professionnelle allait dorénavant être dominé par deux contraintes permanentes : il était propriétaire de Gros Bois et il dirigeait une Étude de notaire que ses prédécesseurs avaient placée dans les premiers rangs des Études Lyonnaises, tout cela alors qu'il n'avait pas encore trente ans. Pour avoir été toute sa jeunesse, comme le souligne le père Rosette, anxieux des responsabilités qui lui incomberaient et de la façon dont il les porterait, il ne pouvait être mieux servi. Dès qu'il fut définitivement chez lui à Gros Bois, il commença à avoir une position officielle dans la commune d'Ouroux. Nommé fabricien par le curé d'Ouroux en 1880, il se trouva honoré d'héberger dès cette même année pendant deux jours le Cardinal Caverot venu à Ouroux pour la Confirmation ; le grand salon de Gros Bois dut être transformé en chapelle pour la messe ; ce fut l'amorce d'une tradition qui, par la suite, vit s'arrêter souvent à Gros Bois les Cardinaux titulaires de l'Archevêché de Lyon lors de leurs tournées en Haut Beaujolais ; Adrien s'y prêta volontiers.

En 1882 il se présenta aux élections municipales à Ouroux et fut élu mais non sans qu'un incident, assez significatif de la mentalité du moment, se produisît. Le Beaujolais qui n'avait pas encore la prospérité que nous lui connaissons aujourd'hui, votait en effet à cette époque radical-socialiste avec une connotation anticléricale et même antireligieuse marquée et celle-ci commençait à gagner le Haut-Beaujolais très attaché jusque là aux valeurs traditionnelles. L'élection d'Adrien, ayant été entachée d'une légère irrégularité de procédure, fut attaquée devant le Conseil de Préfecture. L'Affaire n'eut heureusement pas de suite, mais lorsqu'en 1883, le Conseil Municipal d'Ouroux lui demanda de prendre la Mairie, Adrien savait déjà que tout n'irait pas sans difficultés.

% Page source 14 : IMGP0139.JPG ; Originaux/Page 14.jpg

Enchanté et très fier de cette nomination au poste de premier Édile. Adrien se mit résolument au courant de la gestion de la commune et le père Rosette indique qu'il y consacrait parfois des journées entières, agrémentées de visites au percepteur ou au secrétaire de mairie à Monsols, visites qu'il lui fallait faire quelquefois aller et retour à pied ; l'entraînement acquis avec le cousin Neyrat ne lui était pas inutile.

Des détails de la façon dont il mena la gestion de Gros Bois et notamment de la forêt que F.F.B. avait commencé à constituer 20 ans plus tôt, nous ne savons pas grand chose. Implantée essentiellement sur des versants Nord (Nord-Est à Nord-Ouest), assez pentus, cette propriété se prêtait mal à l'élevage et encore moins à la culture. A.B. acheva de boiser toutes ces pentes en y plantant des résineux (épicéas, sapins, mélèzes). S'occupant très soigneusement de ces plantations, il en fit une forêt qui, avec 140 Ha d'un seul tenant et d'excellents chemins de desserte, devint l'une des plus belles forêts du Haut Beaujolais comme ses petits enfants purent s'en rendre compte quand ils en héritèrent cent ans plus tard.

Par contre il est resté trace de toute une série d'acquisitions de petites parcelles de bois ou de terres voisines de Gros Bois, faites en vue d'arrondir la propriété ou de supprimer des enclaves, ce fut une douzaine de parcelles, 7 Ha environ, qu'Adrien ajouta à son patrimoine entre 1881 et 1890, tandis qu'en contrepartie il vendit en 1882 le petit domaine de Byre, hérité de F.F.B. et qui constituait un écart peu intéressant.

Les années 1880-1890 virent une évolution importante du contexte familial dans lequel le ménage Adrien Berloty-Anne Marie Demoustier les avait abordées.

En 1883, la grand-mère Scholastique Bonaventure Berloty-Poulat s'éteignit à l'âge de 97 ans, elle laissa certainement un grand vide. Comme elle avait gardé jusqu'à sa mort la jouissance du château de Gros Bois et de ses dépendances, elle y recevait en quelque sorte l'épouse de son petit fils et celle-ci ne devait pas s'y trouver pleinement chez elle. Cette situation prit fin avec sa disparition et Anne-Marie fut enfin maîtresse de maison à Gros Bois.

% Page source 15 : IMGP0140.JPG ; Originaux/Page 15.jpg

Elle y faisait des séjours de plus en plus longs en été, s'y trouvant seule avec Marie tandis qu'Adrien restait à Lyon pour l'Étude à moins que sa santé, qui ne le laissait pas tranquille, ne l'obligeât à des cures qui constituaient d'ailleurs un peu pour lui des distractions. En 1884 et en 1885, il fit des saisons à Vichy à la suite de violentes douleurs d'estomac ; on lui fit de pénibles lavages d'estomac mais il y apprit aussi à jouer au billard et puis, satisfaction personnelle certaine, il y changea de médecin, le premier qu'il avait choisi s'étant avéré soutien acharné de l'anticlérical Paul Bert. En 1888, c'est à Aix-les-Bains qu'il passa trois semaines ; ses lettres à Anne Marie décrivent avec humour la clientèle bigarrée, principalement féminine et franco-anglaise, de l'hôtel où il est descendu et il est tout heureux d'être reçu comme un grand homme par le curé local.

Dans l'été 1885, autre événement familial important et celui-ci tout à fait inattendu. Début Août, la charge d'Agent de Change de Romain Demoustier fut déclarée en faillite ; il s'était sans doute quelque peu aventuré. La liquidation de cette faillite se fit sans que Romain fût obligé de vendre l'essentiel des immeubles qu'il possédait à Lyon (la Favorite-place Carnot-Quai de la Charité) mais il dut vendre toutes ses valeurs mobilières et la chose laissa des traces qui subsistaient encore lorsque les successions des deux époux vinrent à être réglées quelques années plus tard. L'inventaire de la succession d'Anne Marie Phelip, l'épouse de Romain, établi en Janvier 1887, fait état en effet d'une créance de 127.000 F. d'Anne Marie sur son époux qui montre probablement qu'il avait fallu puiser dans la fortune d'Anne Marie, mariée sous le régime dotal, pour aider à rembourser tout le passif. Quant à la succession de Romain lui-même décédé en Novembre 1897, elle montre que Romain n'avait pu en 1885 verser à sa fille Elizabeth, qui se mariait en épousant son beau-frère René Brunet-Leconte veuf de sa sœur Marie Christine, la donation en avance d'hoirie de 120.000 F. qu'il lui avait faite par contrat de mariage le 29 Juillet, quelques jours avant la faillite ; c'est seulement en Novembre 1897 qu'Elizabeth put récupérer son dû.

Afin de disposer, en sus de son appartement, place de la Bourse, d'un logement plus agréable pour le printemps ou l'été, époques auxquelles la chaleur à Lyon est souvent très lourde. Adrien et Anne Marie louèrent aux Pères Maristes une villa avec petit jardin située Montée des Carmes Déchaussés. Cette résidence, située au dessus de Saint-Paul, sur la pente Nord de la colline de Fourvière, pouvait effectivement montrer quelque agrément. C'est là que naquit le 15 Août 1886, le fils espéré depuis 10 ans et qui reçut le prénom de Félix ; il complétait heureusement la famille.

% Page source 16 : IMGP0141.JPG ; Originaux/Page 16.jpg

Un mois après le 25 Septembre, Anne Marie Phelip Demoustier décédait dans sa maison de la Favorite. Romain s'y trouvait maintenant tout seul ; il préféra une solution le rapprochant en hiver du centre de Lyon et finalement s'installa place de la Bourse dans le petit appartement du 2ème étage que le décès de Scholastique Bonaventure avait rendu libre ; sa fille Anne Marie Berloty, pouvait s'occuper de lui et il avait deux de ses petits enfants à sa porte ; au printemps il remontait à la Favorite. Très vite il prit l'habitude de faire d'assez longs séjours à Gros Bois en été quand sa fille s'y installait et il y trouvait la possibilité d'y exercer ses talents d'horticulteur, notamment en cultivant des fougères.

Le père Rosette parle longuement de la façon dont A.B. travaillait à l'Étude maintenant qu'il en avait l'entière responsabilité. Toujours arrivé le premier après avoir assisté à la messe à Saint-Nizier ou à Saint-Bonaventure, il l'abandonnait à midi pour le temps du déjeuner et ne la quittait en fin d'après-midi que vers 19 H 30.

Toujours inquiet de ne pas avoir préparé avec assez de soin les affaires qu'il traitait, inquiet aussi des risques du métier en matière financière, domaine dans lequel il reconnaissait lui-même n'être pas très à l'aise, il lui arrivait souvent pour un dossier à traiter le lendemain, de redescendre à l'Étude après dîner ou de se lever à 5 H du matin pour travailler avant d'aller à la messe. Peut-être Anne Marie trouvait-elle alors que la proximité immédiate de l'Étude, juste au dessous de leur appartement, n'était pas sans inconvénients.

À Ouroux les écoles, dont la population était satisfaite, étaient toutes à direction confessionnelle. Lorsque Paul Bert et Jules Ferry préparèrent les premières mesures dirigées contre la liberté de l'enseignement avec comme prétexte de rendre l'enseignement obligatoire, Adrien, alors non encore titulaire de la Mairie, commença à prendre publiquement position contre cette politique et à s'engager ainsi dans une lutte à laquelle il se donna jusqu'à sa mort.

Il avait eu dès 1880 l'occasion de montrer qu'il ne comptait pas se contenter dans ce domaine de protestations platoniques. La loi du 29.3.1880 dispersait les Jésuites et obligeait toutes les congrégations à obtenir dans les six mois l'agrément de l'État. Lorsque cette loi fut publiée Adrien, dont la demande d'agrément comme notaire était en cours d'examen par la Chancellerie, préféra s'abstenir de toute manifestation immédiate mais à partir du 28 Août 1880, il se trouva plus libre pour réagir.

% Page source 17 : IMGP0142.JPG ; Originaux/Page 17.jpg

Ayant appris le 15 Octobre qu'une opération de police était en cours pour chasser les Carmes de leur couvent qu'il connaissait bien puisque voisin de sa résidence de printemps, il quitta l'Étude immédiatement, traversa la Saône vers Saint Paul et s'enferma toute la journée chez les Pères pour les aider à résister à l'opération et pour que les brutalités, inévitables en pareil cas, ne soient pas sans témoin.

Benoît Coste son grand père, qui avait eu des réactions tout à fait comparables au moment où le Premier Empire menait sa politique anticléricale, n'aurait pas fait mieux. Ceci amène à se demander si dès cette époque Adrien ne faisait pas partie de la « Congrégation des Messieurs » cette organisation ultra secrète à la fondation de laquelle Benoît Coste avait participé au début du siècle avec pour objet une lutte active contre les mesures anticléricales et antireligieuses de l'Empire. Adrien ne mentionna jamais cette appartenance à qui que ce soit, les statuts de la Congrégation le lui interdisant de façon absolue ; le père Rosette, qui devait s'en douter, n'en parla pas non plus mais elle est néanmoins des plus probable. Adrien n'aurait en effet, jamais été agréé Membre du Conseil Central de la Propagation de la Foi, comme il le fut par la suite, s'il n'avait pas été au préalable et depuis longtemps membre de cette Congrégation ; il avait dû y adhérer après son mariage en 1875, comme les statuts le prévoyaient.

Peu de temps après c'est à Ouroux et sur un autre plan qu'il se manifesta en défenseur infatigable de l'enseignement libre et des congrégations religieuses. L'école des garçons était en effet gérée par des frères Maristes qui furent mis en demeure de céder la place à un instituteur, conformément aux mesures arrêtées par la IIIème République pour les écoles communales. Les Maristes quittèrent l'école communale et un instituteur fut mis en place. Bien qu'il n'eût rien à reprocher à celui-ci, Adrien réagit en cherchant à maintenir à Ouroux un enseignement religieux par la mise en place d'une école libre à côté de l'école communale. Les démarches entreprises à cet effet n'aboutirent malheureusement pas mais on retrouva là, Adrien adoptant pour ce problème une position activiste tout à fait comparable à celle de son grand père Benoît Coste 60 ans plus tôt.

Devant cet insuccès, Adrien changea de tactique pour l'école des filles. Les religieuses ayant été contraintes à abandonner elles-aussi les locaux qu'elles occupaient pour que l'école publique puisse s'y installer, Adrien décida de faire construire à ses frais, sur deux petites parcelles qu'il avait acquises récemment en plein centre d'Ouroux, des locaux à mettre à la disposition des sœurs pour qu'elles s'y réinstallent. Aucun obstacle n'étant % Page source 18 : IMGP0143.JPG ; Originaux/Page 18.jpg ; raccord avec p. 17
intervenu de la part du Conseil Municipal, tout fut fait rapidement et dès Octobre 1887 les religieuses ouvrirent l'école libre des filles. Ce fut un succès. Comme en outre, pour arriver à ce résultat, Adrien y avait été de sa poche, cela contribua à accroître l'estime dont il commençait à jouir dans les milieux ecclésiastiques de la région Lyonnaise.

La position prise par Adrien dans l'affaire de l'expulsion des Carmes fut de son côté très appréciée des congrégations religieuses installées à Lyon et quelques mois après l'une d'entre elles vint lui demander de l'aider dans la vente d'immeubles lui appartenant. Adrien accepta et c'est ainsi, mais cette fois-ci en tant que notaire et non plus à titre strictement individuel, qu'Adrien devint le conseil de diverses congrégations pour des opérations analogues, d'ailleurs de plus en plus fréquentes.

Cela lui valut d'être invité à participer à de nombreuses réunions dans la région et même dans tout le Sud-Est et à discuter sur la stratégie de défense à adopter. La loi du 29.12.1884 ayant donné un tour de vis supplémentaire sur les biens des congrégations. Adrien jugea qu'il était préférable d'adopter en face de cette loi une attitude passive plutôt que de se ruiner pour se maintenir en acceptant les paiements exorbitants qu'elle prévoyait. Avec son beau-frère, Joseph Rambaud, il développa ces idées dans une brochure intitulée : « Ce qu'il faut faire face aux lois Brisson ».

Cette brochure eut un certain succès mais néanmoins limité et plusieurs congrégations ne suivirent pas ses conseils. Adrien en ressentit quelque amertume et même des lettres que cite le père Rosette montrent l'apparition d'une crise de découragement analogue à celles qu'il avait connues autrefois.

D'autres lettres, quelque peu postérieures, que cite le père Rosette, montrent également qu'Adrien ne fut pas sans inquiétude devant la politique générale suivie alors par les autorités ecclésiastiques. Ces lettres furent écrites en 1889 de Rome où il était allé, à titre de membre de la Commission des Finances de la Propagation de la Foi, pour rencontrer les services pontificaux responsables des missions à l'étranger. Les nombreux prélats qu'Adrien vit à cette occasion sont classés par lui en deux catégories : les fidèles de Pie IX, pape intégriste s'il en fut, qui ont droit à toute sa sympathie et les libéraux, fidèles de Léon XIII, chez lesquels il n'accepta à dîner que contraint et forcé ; c'était là aussi une position tranchée du style Benoît Coste. Il se rendit néanmoins, il ne pouvait faire autrement, à la réception que Léon XIII accorda à la délégation de la Propagation de la Foi, et d'ailleurs, très honnêtement, il fut très sensible à l'accueil que Léon XIII lui réserva. Ce fut Léon XIII qui, quelques années plus tard, le nomma Chevalier de Saint Grégoire le Grand.

% Page source 19 : IMGP0144.JPG ; Originaux/Page 19.jpg

Quelques mois après la naissance de Félix, probablement en 1887, Adrien abandonna la résidence secondaire qu'il louait depuis longtemps Montée des Carmes Déchaussés pour en louer une autre rue du Juge de Paix à côté de la basilique de Fourvière. La maison possédait une terrasse de laquelle par temps clair on voyait très bien les Alpes du Dauphiné et le Vercors. Ce n'était plus à deux pas de l'Étude mais par contre Anne Marie était plus près de la Favorite où elle pouvait aller facilement voir son père et emmener ses enfants.

En dépit des apparences et de la façon dont Adrien se comportait à Gros Bois et à Ouroux, il n'y était pas encore pleinement chez lui, l'indivision établie avec son frère Bonaventure au 1er Janvier 1879 subsistant toujours. Celui-ci âgé de 35 ans, avait maintenant terminé son noviciat, le maintien de cette indivision ne se justifiait plus et en Octobre 1890 les deux frères réglèrent le problème, Adrien rachetant à Bonaventure tous ses droits, en contrepartie principalement des avances qu'il lui avait faites depuis son entrée au noviciat.

Février 1891, Adrien a maintenant quarante ans, une nouvelle naissance est attendue par le ménage et une seconde fille Elizabeth, viendra en Septembre rejoindre Marie qui a 12 ans et Félix qui en a 5, avant que viennent s'ajouter Joseph en 1893 et Paul en 1895.

Malgré ses nombreuses occupations, Adrien allait consacrer beaucoup d'attention à l'éducation de ses enfants et spécialement, au début, à celle de Marie et de Félix. Les lettres, citées par le père Rosette, qu'il leur écrivait lui-même de Lyon quand en été, ils séjournaient à Gros Bois, étaient toujours très gentiment affectueuses mais ne craignaient pas de manifester, quand il le fallait une certaine sévérité ; c'est par elles que l'on apprend que Marie, qui venait d'être envoyée en Angleterre, était fermement invitée à se surveiller car, à 12 ans, elle se montrait altière avec les domestiques, manquant trop souvent d'indulgence voire même de charité.

Sa santé continuait à lui procurer quelques soucis et il dut faire une nouvelle saison à Vichy. Il se montrait en même temps très attentif à la santé de ses enfants. Il gardait mauvais souvenir de ses bronchites aussi quand en 1890, Félix en attrapa une assez violente, Adrien n'hésita pas à envoyer Anne Marie et ses enfants faire un séjour à Montreux pour l'en remettre complètement. De même en 1894 et 1895 séjour à la Bourboule et à Uriage pour soigner à nouveau la bronchite de Félix ; c'est Anne Marie qui y va, Adrien ne fait qu'y passer.

% Page source 20 : IMGP0145.JPG ; Originaux/Page 20.jpg

Malgré le peu de succès des placements immobiliers que F.F. Berloty avait faits à l'étranger notamment à Cadix, Adrien est comme son père attiré par des placements de cette nature. En 1891 avec ses amis Ernest Finaz et Lucien Souchon, il achète en plein centre de Tunis, à l'angle de la rue d'Italie et de la rue d'Angleterre, un immeuble de deux étages sur rez-de-chaussée appelé « Maison de la Poste ». L'achat fut payé 220.000 F. dont un cinquième par Adrien, le solde par Finaz et Souchon. L'immeuble était bien placé et l'achat s'avéra intéressant car en 1905 au décès d'Adrien il fut estimé à 330.000 F. pour un loyer annuel de 17.500 F. Cet achat amorça une série de voyages en Tunisie et une autre acquisition qu'Adrien fit par la suite avec ses amis en achetant début 1899, un terrain non construit d'une superficie de 2 Hectares à Sfax.

Il était toujours partisan de l'école libre et farouchement hostile aux mesures antireligieuses et anticléricales de la IIIème République.

Plusieurs des lettres citées par le père Rosette regrettent franchement la position trop conciliante que prenaient sur ce point certaines personnalités importantes des milieux catholiques parisiens, le Sud-Est de la France restant en général plus intégriste ; Adrien avait pourtant accueilli favorablement l'Encyclique Rerum novarum publiée par Léon XIII en 1891. À vrai dire si le père Rosette est prolixe sur tout ce qui concerne l'action menée par Adrien contre les mesures antireligieuses, il est beaucoup plus discret sur ce que pouvaient être ses pensées intimes face à l'évolution des mentalités, et par exemple il ne dit pas un mot de l'affaire Dreyfus qui secoua la France dans les années 1894-1898 avec des répercussions politiques considérables devant laquelle Adrien ne resta certainement pas indifférent.

La loi de Décembre 1884 concernant les congrégations s'étant avérée d'application difficile, une nouvelle loi fut votée en Avril 1895, moins odieusement contraignante sur certains points, mais plus exorbitante sur le plan fiscal et financier, Adrien réagit comme précédemment en publiant avec son beau-frère Rambaud une note « Étude pratique sur la situation faite aux Congrégations par la loi fiscale de 1895 » note dans laquelle, entre autres mesures à prendre, il prônait à nouveau la résistance passive.

Pendant tout l'été 1895, il s'employa à diffuser cette note en la commentant notamment dans les milieux catholiques du Sud-Est où elle fut en général bien accueillie, alors que le milieu parisien plus mêlé aux négociations avec le pouvoir était plus partisan de solutions de compromis. Il renonça à ses vacances à Gros Bois, ne pouvant y travailler correctement au problème du moment.

% Page source 21 : IMGP0146.JPG ; Originaux/Page 21.jpg

Ce fut certainement pour lui une période d'intense activité marquée comme le signale le père Rosette par de nombreux voyages à Paris et en province. Très déçu par la division qui se manifestait dans la presse catholique française partagée à ce moment entre la tendance intégriste de Louis Veuillot et la tendance libérale dérivée de Montalembert et de Rerum novarum, il lui arriva de se laisser aller certains jours à des crises de découragement qui heureusement ne duraient pas. Il fut d'ailleurs certainement conforté dans sa position lorsqu'il fut au cours de cette même fin d'année 1895 agréé comme membre du « Conseil de Lyon » de la Propagation de la Foi, fonction à laquelle il avait été bien préparé par sa participation depuis de nombreuses années à la Commission des Finances de cette organisation.

Rien ne témoigne plus de la vigueur de ses sentiments à cette époque, que la réaction publique qu'il eut en Mai 1897 comme suite à un incident politique consécutif à l'incendie du Bazar de la Charité à Paris. Cet incendie, qui fit plus de 200 morts, émut l'opinion et il y eut des funérailles officielles à Notre Dame de Paris ; au cours de celles-ci le prédicateur jugea opportun de parler des châtiments que la Providence réserve aux peuples et aux nations qui le méritent. Le lendemain 18 Mai, le Président du Conseil prenant spécialement la parole à la Chambre des Députés releva la chose dans un discours violemment antireligieux, discours dont immédiatement la Chambre vota l'affichage obligatoire dans toutes les mairies de France. Adrien savait ce qu'il avait à faire et, le 22 Mai, il informait par lettre le Préfet du Rhône qu'à Ouroux le discours antireligieux ne serait pas affiché, il signait sa lettre « Adrien Berloty, \textbf{maire de Saint-Antoine d'Ouroux} ». Il était bien le digne petit fils de Benoît Coste ; suivirent de nombreuses lettres de félicitations venant de toute la France et même de l'étranger. Quant au Préfet, qui ne pouvait guère faire autrement, il suspendit le Maire d'Ouroux et Adrien se retrouva simple Conseiller Municipal.

Une vie d'homme public, telle que la menait Adrien, est ainsi souvent parsemée d'événements inattendus ; Adrien venait de s'en apercevoir. Elle lui avait aussi procuré auparavant l'occasion d'intervenir lors de l'assassinat du Président Carnot à Lyon en Juin 1894. Le banquet officiel avait eu lieu dans le palais de la Bourse à côté de l'Étude ; Adrien y avait été invité en qualité de Maire d'Ouroux, il venait de rentrer à l'Étude après la cérémonie, quand il vit par la fenêtre l'assassin se précipiter sur le Président. Aussitôt descendu, il constata que dans la confusion générale il n'y avait pas de prêtre. Il courut à Saint Jean prévenir le Cardinal Coullié qui arriva à temps à la Préfecture pour administrer le Président.

% Page source 22 : IMGP0147.JPG ; Originaux/Page 22.jpg

Tout cela ajouté au travail courant de l'Étude l'occupait beaucoup et le père Rosette cite de lui des lettres dans lesquelles il se dit épuisé après une journée commencée à 4 H du matin. Mais cela ne l'empêchait pas de s'occuper activement de sa famille, de ses résidences et de sa fortune personnelle.

En Juillet 1896, pour accélérer le rétablissement de Félix qui venait d'être opéré des végétations, opération encore assez rare à l'époque, et pour en même temps montrer à ses enfants autre chose que Gros Bois, Adrien loua une villa à Arromanches à quelques kilomètres de Caen. Anne-Marie y arriva la première avec les enfants ce qui leur permit de découvrir Paris au passage ; Adrien les y rejoignit pour quelques jours de vraies vacances loin de l'Étude et de la Mairie.

L'année d'après pour permettre à Félix de récupérer après une année à Mongré qu'il a mal supportée, c'est au Val André que la famille se rend en Juillet et Félix s'en trouva bien pour rentrer en Octobre chez les Jésuites de la rue de Sèze à Lyon.

Au même moment, pour une raison que l'on ignore A.B. abandonna au printemps sa résidence secondaire de la rue du Juge de Paix pour louer une villa avec jardin située Chemin de St.~Irénée à Sainte Foy à hauteur de l'hôpital Debrousse. Cette villa avait elle aussi certainement une belle vue sur les Alpes mais elle était sensiblement plus loin de l'Étude et, quand il s'y installa au printemps, Adrien fut obligé de faire venir de Gros Bois voiture, cheval et cocher pour le conduire chaque jour au funiculaire de Saint-Just par lequel il descendait facilement à la Bourse.

Tous ces mouvements montrent qu'Adrien n'était pas d'un naturel casanier, bien au contraire.

Le 1er Novembre 1897, Romain Demoustier décéda dans sa maison de la Favorite. Depuis la mort de sa femme, Anne-Marie Phelip, il s'appuyait beaucoup sur sa fille Anne Marie Berloty et, comme il occupait en hiver le petit appartement du 2ème étage de la place de la Bourse, la famille d'Adrien le voyait beaucoup. Dans ses « Souvenirs d'enfance » Félix parle longuement des visites très fréquentes qu'il faisait à la Favorite et dont il garde un excellent souvenir.

Le décès de Romain fut certainement sensible à toute la famille aussi en raison de ce deuil Adrien et Anne-Marie souhaitèrent ne pas avoir à participer aux festivités, familiales ou autres, traditionnelles en fin d'année.

% Page source 23 : IMGP0148.JPG ; Originaux/Page 23.jpg

Ils emmenèrent leurs enfants quelques jours à Tamaris à côté de Toulon entre Noël et le Jour de l'An et ce fut pour eux l'occasion de faire connaissance avec la Méditerranée, bien que la période n'y soit pas très favorable. Vraiment, pour la famille d'Adrien, les voyages faisaient partie de la vie courante.

Le petit appartement du 2ème étage de la place de la Bourse étant libéré, maintenant que Romain Demoustier n'y était plus, Adrien le récupéra pour son usage personnel afin d'agrandir son logement devenu un peu exigu depuis que ses enfants avaient grandi. Il en profita pour redistribuer les pièces et aussi améliorer le confort, en y installant l'éclairage électrique ; encore très peu utilisé à cette époque ; l'éclairage se faisait avec des lampes à filaments de carbone peu lumineuses mais c'était un progrès considérable par rapport aux lampes à pétrole qui exigeaient un entretien quotidien. Il y installa aussi, innovation importante, des W.C. à chasse d'eau ; par contre on en resta à un seul cabinet de toilette pour les parents et les enfants et toujours pas de salle de bains ; il fallait aller prendre un bain dans un établissement de bains publics.

Dans l'appartement ainsi agrandi, Adrien et Anne-Marie furent plus à l'aise pour recevoir, et Félix Berloty dans ses « Souvenirs d'enfance » décrit en détail les mondanités qu'il voyait se dérouler chez ses parents à cette époque. Ceux-ci avaient à leur service quatre domestiques : cuisinière, femme de chambre, bonne d'enfants et un valet de chambre qui couchait d'ailleurs à l'Étude où il avait la charge de l'entretien de toutes les lampes à pétrole. Ils disposaient pour recevoir d'un très grand salon et d'une salle à manger de 30 couverts. À déjeuner Adrien accueillait souvent des clients mais aussi quelques amis fidèles. Par exemple l'abbé Boisard qui a laissé un nom à Lyon, ayant été fondateur de la première école d'apprentissage pour les métiers du bois et des métaux, ou encore Stanislas Neyrat, son cousin de voyage en Espagne, devenu avec l'âge protonotaire apostolique et doyen du chapitre de Saint-Jean.

Si tous les jeudis soirs, jusqu'au décès de sa belle-mère Anne-Marie Demoustier Phelip, Adrien emmenait ses enfants dîner chez leurs grands parents, cette réception rituelle fut remplacée par la suite par des dîners avec les Boffard et les Rambaud, alternativement chez les uns et chez les autres. Il y avait aussi quatre ou cinq fois par an, place de la Bourse à 19 heures, dîner de gala en habit et robe du soir pour 25 ou 30 couverts, suivi d'une réception plus nombreuse qui se terminait à minuit, avec quelquefois de la musique, Anne-Marie Demoustier tenant le piano, Adrien et sa fille Marie chantant.

% Page source 24 : IMGP0149.JPG ; Originaux/Page 24.jpg

Le dimanche Adrien allait souvent au concert, en général à la Scala en y emmenant ses aînés, ou bien assistait à des conférences ou des réunions politiques ou religieuses. Il n'allait que très rarement au théâtre ou à l'opéra de Lyon, et à Paris ne fréquentait guère que l'Opéra ou la Comédie Française. Le père Rosette prend bien soin de souligner qu'il s'abstenait de spectacles lestes ou déshabillés.

En hiver il y avait aussi pour Félix et Marie des leçons d'équitation au manège Pochet rue Molière, et Adrien qui ne semble pas avoir beaucoup pratiqué le cheval jusque là, puisque le père Rosette mentionne que c'était à pied qu'il allait de Gros Bois à Monsols pour le service de la Mairie, se mit à suivre également de temps en temps ces leçons. Bien que touchant à la cinquantaine et malgré sa santé déjà défaillante, il lui arrivait ainsi de retrouver de temps en temps de véritables moments de jeunesse.

À l'approche des années 1900 qui, comme on le verra ci-dessous, apporteront un nouveau changement profond dans la vie d'Adrien, son souci d'arrondir Gros Bois, chaque fois que l'occasion s'en présentait, restait toujours aussi vif. Après avoir acheté en Janvier 1897, un hectare et demi de très bons prés de rivière à la Tallebarde en bordure de la Grosne, à côté d'Ouroux, il acheta en Mai 1898, les 25 Hectares supplémentaires de terres et de prés du domaine des Friands situé de l'autre côté d'Ouroux sur le versant Nord de la vallée de la Grosne et donc non contigu à Gros Bois.

Son intérêt pour le notariat ne se ralentissait pas non plus, car c'est à ce moment qu'il accepta de devenir en Janvier 1900, Président de la Chambre des Notaires après avoir été depuis 1890 membre du Conseil de cette Chambre et notamment syndic pendant deux années. C'était endosser de nouvelles responsabilités, en même temps que cette désignation lui apportait le témoignage de l'estime de ses collègues ; il prit officiellement ses fonctions de Président en Janvier 1900.

\begin{center}\rule{0.5\linewidth}{0.5pt}\end{center}

% Page source 25 : IMGP0150.JPG ; Originaux/Page 25.jpg

\chapter{La retraite et les Massues}\label{la-retraite-et-les-massues}

L'année 1900 commença sans que rien ne fît prévoir une évolution brutale de l'activité d'Adrien. Au cours de l'été Adrien et Anne-Marie emmenèrent leurs deux aînés, Félix et Marie, accompagnés de leur cousine Marguerite Brunet-Leconte, dans un long voyage en Bavière et en Autriche ; départ par la Suisse, Zurich, les chutes du Rhin à Neuhausen puis Munich avec une représentation de la Passion à Oberammergau, ensuite Salzbourg et le lac Königssee, trouvé absolument enchanteur, trois jours à Vienne retour par Innsbruck et l'Arlberg ; une vraiment bonne détente avant d'aborder l'hiver 1900 - 1901 au cours duquel Félix ferait sa troisième, et pourtant d'importants changements étaient certainement déjà envisagés à ce moment avec l'achat de la propriété des Massues et l'abandon de l'Étude, décisions sans doute liées l'une à l'autre et qui durent surprendre la famille aussi bien que les milieux professionnels, car Adrien abandonnait son Étude sans même avoir atteint cinquante ans.

L'acte d'achat des Massues fut signé le 30 Octobre 1900 et c'est dans les semaines qui suivirent qu'Adrien annonça qu'il souhaitait céder l'Étude au cours de l'été 1901. Aucun document ne permet de connaître avec précision les motifs de ce bouleversement ; on peut en supputer plusieurs.

À cette époque, Adrien supportait de plus en plus difficilement les fatigues résultant des activités multiples qu'il assumait. Sa santé déclinait au point que, à en juger d'après la lecture de certaines de ses lettres, il semble avoir été repris par de véritables crises de découragement devant le risque d'être obligé de s'arrêter impromptu. En ce cas il lui faudrait céder l'Étude sans avoir la certitude de la voir revenir plus tard à son fils. Si en effet Félix, né en 1886, en venait à souhaiter prendre la suite de son père plutôt que de s'orienter vers une autre activité, ce serait seulement vers les années 1915 à 1920 que la chose pourrait intervenir, Félix ayant alors à ce moment, avec la trentaine, la possibilité de demander avec succès son agrément comme notaire. Au cas où d'ici là, Adrien serait conduit à s'arrêter, il lui faudrait trouver un successeur qui accepte de ne prendre l'Étude que pour un nombre d'années limité et ce ne serait pas facile. Alors, puisqu'une telle éventualité paraissait inévitable compte tenu de sa santé, ne valait-il pas mieux réaliser de suite la cession en prenant le temps voulu pour trouver un successeur qui accepte de réserver la position de Félix pour l'avenir.

Il se peut aussi qu'afin d'être plus libre de sa position dans la défense des congrégations religieuses, Adrien ait souhaité être débarrassé des contraintes du notariat et de la présidence de la Chambre des Notaires, l'arrivée au pouvoir du Ministre Combes, faisant prévoir une reprise immédiate et vigoureuse de la lutte antireligieuse.

% Page source 26 : IMGP0151.JPG ; Originaux/Page 26.jpg

Si l'Étude était ainsi cédée à un tiers, avoir sa résidence contiguë à celle-ci, place de la Bourse ne présentait plus d'intérêt et même il valait mieux s'en éloigner pour que tous les liens soient effectivement bien rompus. L'achat des Massues tombait à point pour apporter à ce problème une excellente solution, le ménage d'Adrien abandonnant la résidence principale dans laquelle il vivait depuis 25 ans pour la transférer dans une propriété agréable, de semi-campagne quoique toujours proche du centre de Lyon ; Adrien y trouverait un environnement plus favorable pour sa santé.

Située 60 Chemin des Massues à la limite Ouest du 5ème arrondissement de Lyon, la propriété avait une superficie importante pour une localisation en très proche banlieue, 2 Hectares 80 avec parc, pelouse, jardin d'agrément et jardin potager, mais avec des bâtiments en très mauvais état qu'il faudrait démolir et remplacer. Cette propriété appartenait alors à Monsieur Desprez qui, retiré à Nice, n'y résidait pratiquement plus ; elle était voisine de la propriété Phelip et, de ce fait, déjà bien connue d'Anne-Marie et d'Adrien.

Elle fut achetée 72.500 F, par Anne-Marie et non par Adrien. Peut-être celui-ci, ayant le pressentiment qu'il ne ferait pas de vieux jours, choisit-il cette solution pour assurer à son épouse, s'il venait à disparaître le premier, une propriété bien à elle alors qu'à Gros Bois elle ne serait jamais que chez ses enfants.

Pour l'Étude Adrien trouva rapidement un successeur possible en la personne de Louis Charrat, un ancien condisciple de Mongré, qui était déjà notaire à Belleville. Celui-ci s'engagea à céder l'Étude à Félix, lorsqu'il prendrait sa retraite, et au prix auquel il l'achetait à Adrien. Tout fut ainsi traité au mieux.

Corrélativement Adrien résilia dans l'été 1901 ses fonctions de Président de la Chambre des Notaires, ce qui fut l'occasion d'une manifestation de grande sympathie de tous les notaires Lyonnais, comme cela ressort des discours prononcés et dont le père Rosette cite de nombreux passages. Son successeur à la présidence, Maître Chaine, proposa aussitôt à ses collègues de demander au Ministre de la Justice qu'Adrien soit nommé « Notaire Honoraire », distinction qui n'était pas accordée systématiquement.

L'été à Gros Bois fut agité par de nombreuses visites préparatoires à tous ces changements jusqu'au jour où, Louis Charrat ayant reçu l'agrément du Ministre de la Justice le 7 Septembre 1901, acheta immédiatement l'Étude, cet achat fut fait au prix de 500 000 F, Adrien restant propriétaire des locaux.

% Page source 27 : IMGP0152.JPG ; Originaux/Page 27.jpg

Maintenant à la retraite, Adrien eut tout le loisir de s'occuper des travaux de construction de la nouvelle villa des Massues. Avec l'aide de l'architecte choisi à cet effet, Monsieur Bonivant, il voulut avoir là une résidence dotée du confort le plus moderne à l'époque avec bien évidemment l'électricité mais aussi avec cabinet de toilette dans presque toutes les chambres et salle de bains alimentée en eau chaude par une bouilloire spéciale du fourneau de la cuisine située à proximité immédiate.

À Pâques 1902, il s'installa aux Massues et tout de suite le ménage noua des relations amicales avec le voisinage avec les Mouterde, avec Mme de Prandières et Mme de Bonnevie auxquelles A.B. avait d'ailleurs rendu service précédemment ; c'étaient de fréquents déjeuners ou goûters chez les uns et chez les autres.

Autre grande nouveauté pour l'époque, Adrien venait d'acheter une voiture automobile, étant d'ailleurs l'un des tous premiers Lyonnais à faire une telle acquisition. L'été précédent, la Société Audibert et Lavirotte avait gracieusement mis à sa disposition une voiture de sa fabrication. C'est donc une Audibert et Lavirotte, 12 chevaux, 9 Places qu'il acheta ; il disposa bien entendu d'un chauffeur mais il prit son permis de conduire témoignant ainsi d'une jeunesse d'esprit incontestable. Il faisait souvent conduire Félix à sa place, et celui-ci, grâce à une faveur spéciale obtint son permis le 8 Août 1902 bien qu'il n'eût pas encore les 18 ans réglementaires. Cette voiture servit d'abord principalement pour les voyages à Gros Bois ce qui, vu l'état encore un peu primitif du matériel et celui non moins primitif des routes, n'allait pas sans problèmes ; il était rare qu'un trajet ne s'accompagne pas de crevaisons. Cette première voiture, construite comme un char à bancs équipé d'un moteur, était trop grande avec ses neuf places et trop lourde. Adrien la remplaça rapidement par une Torpédo Berliet, 4 places plus maniable et, sur le plan mécanique, en progrès très appréciable par rapport à la précédente.

La lutte antireligieuse ayant été relancée avec vigueur par la loi Combes de 1901, Adrien qui, comme indiqué ci-dessus, avait dû sentir venir la chose, reprit en quelque sorte du service en aidant les congrégations en difficulté ou en protestant contre certaines mesures particulièrement vexatoires.

En fin 1901, au titre de trésorier de l'œuvre des Écoles Catholiques de Lyon, il signa une protestation véhémente contre des mesures touchant cette œuvre. Le Préfet du Rhône prit prétexte de la chose pour refuser sa nomination « Notaire Honoraire », refus sur lequel il accepta finalement de revenir.

% Page source 28 : IMGP0153.JPG ; Originaux/Page 28.jpg

On le vit aussi en Octobre 1902 prendre la plume à la demande de la Chambre des Notaires pour rédiger une étude faisant ressortir le non sens d'une proposition de loi de Clemenceau tendant à la fonctionnarisation des notaires.

Ce sont cependant surtout les congrégations qui lui redonnèrent de l'activité. Pendant deux ans, elles afflueront chez lui pour lui demander des conseils mais aussi pour lui demander une aide très spécifique dans l'achat de propriétés en France et à l'Étranger. Cela le conduisit à faire de 1901 à 1903 de nombreux voyages d'abord inévitablement à Paris mais aussi en Angleterre, à Jersey, en Belgique, en Italie.

Si ces voyages souvent rapides le fatiguaient maintenant, il semble qu'il ne s'en souciait pas outre mesure. Il ne savait pas résister au plaisir des voyages, le père Rosette lui faisait dire : « qu'il aimait rouler », et d'ailleurs il ne recula pas à ce moment devant deux voyages assez longs en sus de ceux que lui suscitaient les congrégations, l'un en Suisse en 1901 l'autre à Tunis en 1902, en Suisse, ce fut pour renouveler, en passant par Lucerne, le pèlerinage à Einsiedeln qui avait tant séduit Benoît Coste.

En Octobre 1903, donc à la veille du premier accident de santé sérieux qui sera évoqué ci-dessous, c'est un nouveau voyage à Rome. La Propagation de la Foi vient en effet de décider l'envoi d'une délégation pour présenter au Pape Pie X, élu récemment, les hommages des deux Conseils Centraux de Lyon et de Paris et ceux de l'Œuvre tout entière. Les Présidents des Conseils de Lyon et de Paris n'ayant pu participer à cette délégation qui comportait six personnes, ce fut Adrien qui fut désigné pour la présider ; Membre du Conseil de Lyon depuis 1895, il se trouvait être le plus ancien. Il eut donc l'honneur de s'adresser personnellement au Pape et fut rempli d'aise par la réponse de Pie X qui le conforta dans les idées qu'il avait toujours eues en face de la politique antireligieuse de la IIIème République.

\chapter{Les dernières années}\label{les-derniuxe8res-annuxe9es}

Le 8 Décembre 1903, c'était le pèlerinage traditionnel des Lyonnais à Fourvière. Bien qu'il souffrît de plus en plus d'essoufflements qui le fatiguaient beaucoup, Adrien, soucieux de faire vraiment un pèlerinage, décida d'y aller à pied depuis les Massues. C'était une marche de près de 3 kilomètres avec une pente très raide entre Trion et la porte de Loyasse. Il eut toutes les peines du monde à gravir cette pente et finalement dut s'arrêter avant Fourvière sans avoir rejoint le cortège principal. Le coup était dur.

% Page source 29 : IMGP0154.JPG ; Originaux/Page 29.jpg

Malgré cet avertissement sévère, Adrien ne changea pourtant guère ses habitudes évitant seulement les montées trop fatigantes.

Un mois plus tard dans la nuit du 17 Janvier 1904, une crise de suffocation très violente le saisit. Il dut se mettre à un régime sévère et fut obligé de réduire sérieusement ses activités sans toutefois s'arrêter complètement de travailler.

Après plusieurs semaines de soins, son état s'améliora et au printemps 1904, à la demande de ses amis d'Ouroux il accepta, bien qu'à contre cœur de se présenter au Conseil Municipal dont il avait démissionné en 1900 lors de sa nomination à la présidence de la Chambre des Notaires, il fut élu mais seulement au second tour, ce qui montre que, même à Ouroux, les mentalités avaient évolué.

Passant tout l'été 1904 à Gros Bois pour se soigner et y restant jusqu'en Décembre, il put effectivement participer régulièrement aux séances du Conseil Municipal. L'hiver 1905 ne se passa pas trop mal mais A.B. dut cesser toute activité. Il mit un point final à son testament le 21 Avril ; il se préparait à la mort, les lettres que cite le père Rosette sont éloquentes à cet égard. Il ne pouvait plus se mouvoir sans difficulté et fit installer son bureau au rez-de-chaussée des Massues. Les grandes chaleurs de Juillet le décidèrent cependant à monter à Gros Bois dans les premiers jours d'Août.

Le père Rosette décrit avec émotion les derniers mois de sa vie, ses souffrances, ses nuits d'insomnie, la consolation de pouvoir chaque jour entendre la messe que, par autorisation spéciale de l'Archevêché, on célébrait dans le grand salon de Gros Bois.

La fin approchait, il fallait redescendre à Lyon et Adrien fit ses adieux à Gros Bois le 21 Septembre. La montée à Avenas, dans l'automobile conduite par Félix, amena une oppression angoissante et fit craindre le pire ; une fois dans la plaine de la Saône, cela alla un peu mieux mais la nuit aux Massues fut très mauvaise et Adrien fut administré. Les jours suivants, un léger mieux permit quelques visites et notamment celle du Cardinal Coullié, Archevêque de Lyon, venu personnellement lui témoigner de l'estime qu'il avait acquise auprès du Clergé Lyonnais. Bonaventure fut appelé d'urgence d'Angleterre et resta jusqu'au 8 Octobre. Le départ de Félix au service militaire à Bourg créa une grosse émotion mais celui-ci put revenir le voir vers le 15 Octobre. Il y eut un léger mieux pendant un mois puis l'enflure des jambes reprit de plus belle et Adrien s'éteignit le 15 Décembre 1905 sans achever sa 55ème année.

% Page source 30 : IMGP0155.JPG ; Originaux/Page 30.jpg

\chapter{Portrait et succession}\label{portrait-et-succession}

Dans ses « Souvenirs d'enfance » et peut-être pour corriger l'impression qui ressort de l'ouvrage du père Rosette, Félix dit de son père :

« Malgré sa foi profonde, il restait très humain dans sa vie quotidienne, bon, gai, entrain sans être bruyant, s'intéressant à toutes les manifestations de la vie et à la gestion de sa fortune pour en assurer la conservation et faire plaisir à ses enfants. Il riait et plaisantait volontiers. Il savait ne pas être accaparé par ses occupations et se distraire en restant gai ».

On a vu à quel point il aimait voyager et quelle jeunesse d'esprit il sut conserver jusqu'à ses dernières années. En fait sur le plan religieux, spirituel et probablement aussi politique, il était très proche de Benoît Coste, son grand-père. Sur le plan de la vie familiale et de la gestion de la fortune, il tenait plutôt de François Félix, dont il avait hérité son attrait presque immodéré pour les propriétés rurales.

\begin{center}\rule{0.5\linewidth}{0.5pt}\end{center}

Son testament assez long, comportait d'abord trois pages de considérations à caractère notarial pour préciser comment il voyait les reprises et récompenses de sa succession en ce qui concernait son épouse et ses frères et sœurs. Il léguait à son épouse Anne-Marie jouissance de la maison de Gros Bois avec mobilier et dépendances, pour qu'elle puisse y recevoir ses enfants, à charge pour elle de payer les impôts, les assurances et certains gages. Pour la conservation de la propriété il préférait ne donner aucune instruction laissant à ses enfants et son épouse le soin de choisir le meilleur parti sans se laisser entraîner par d'uniques sentiments personnels ni par les souvenirs de famille.

Il laissait une succession importante, assez complexe à cause des nombreux biens immobiliers qu'elle comportait, mais avec un passif insignifiant. Elle se régla par un partage en date du 28.05.1907. Seule Marie était majeure, tandis que Félix avait été dès le décès de son père, émancipé avec comme curateur, J.P. Boffard, Elizabeth, Joseph et Paul, tous trois mineurs étant représentés par leur subrogé tuteur, J.P. Boffard également.

% Page source 31 : IMGP0156.JPG ; Originaux/Page 31.jpg

La propriété de Gros Bois fut laissée provisoirement en indivision, Anne-Marie Demoustier en conserva l'usufruit conformément au testament de A.B. ; elle recevait en pleine propriété les Escouannes et le Siber ainsi que des valeurs mobilières tandis que chacun des cinq enfants recevait en pleine ou nue propriété le cinquième du solde des valeurs mobilières. Les terrains situés à Saint-Étienne et qui, provenaient de la succession de Joséphine Coste, étaient en indivision entre A.B. et ses cousins, François Coste et Madame de Rochetaillée, furent mis en vente. Le partage définitif n'intervint que longtemps après par un acte sous-seing privé de 1917 et par l'acte du 6 Janvier 1932 liquidant la succession d'Anne Marie Berloty Demoustier décédée le 17 Avril 1931.

\begin{center}\rule{0.5\linewidth}{0.5pt}\end{center}

Le 1er Mai 1993

Léon Peillon
\end{document}
